% META: {"id": "1SPE-SUITES-FR-R2", "chapitre": "1SPE-SUITES", "type_objet": "remediation", "capacites_codes": ["R2"], "status": "generated"}

%---------------------------------------------------------------
% FICHE DE REMISE À NIVEAU — R2 : Pourcentages
% Coefficient multiplicateur, évolutions successives
%---------------------------------------------------------------

\subsection*{Rappel — Pourcentages et coefficient multiplicateur}

\textbf{Coefficient multiplicateur (CM) :}

Appliquer une évolution de $t\,\%$ revient à multiplier par le coefficient multiplicateur :
\[
  CM = 1 + \frac{t}{100} \quad \text{(hausse)} \qquad CM = 1 - \frac{t}{100} \quad \text{(baisse)}.
\]

\textbf{Évolutions successives :} les coefficients multiplicateurs se \emph{multiplient}.

Si on applique successivement des évolutions de CM $q_1$, $q_2$, \ldots, $q_k$, le coefficient global est $q_1 \times q_2 \times \cdots \times q_k$.

\textbf{Lien avec les suites géométriques :} si une grandeur évolue de $t\,\%$ à chaque étape, elle suit une suite géométrique de raison $CM = 1 + \dfrac{t}{100}$.

%---------------------------------------------------------------

\begin{exercice}{1SPE-SUITES-FR-R2-EX1}{1}{5}
Pour chaque situation, donner le coefficient multiplicateur correspondant.
\begin{enumerate}
  \item Une augmentation de $8\,\%$.
  \item Une réduction de $15\,\%$.
  \item Une augmentation de $0{,}5\,\%$.
  \item Une réduction de $100\,\%$.
\end{enumerate}
\end{exercice}

% BEGIN-VERIFY
% # Verification des coefficients multiplicateurs
% CM1 = 1 + 8/100
% assert abs(CM1 - 1.08) < 1e-9
% CM2 = 1 - 15/100
% assert abs(CM2 - 0.85) < 1e-9
% CM3 = 1 + 0.5/100
% assert abs(CM3 - 1.005) < 1e-9
% CM4 = 1 - 100/100
% assert abs(CM4 - 0) < 1e-9
% END-VERIFY

\begin{corrige}{1SPE-SUITES-FR-R2-EX1}

\commentaireMarge{Une hausse de $t\,\%$ donne $CM = 1 + t/100$.}

\textbf{Question 1.} Augmentation de $8\,\%$ : $CM = 1 + \dfrac{8}{100} = 1{,}08$.

\commentaireMarge{Une baisse de $t\,\%$ donne $CM = 1 - t/100$.}

\textbf{Question 2.} Réduction de $15\,\%$ : $CM = 1 - \dfrac{15}{100} = 0{,}85$.

\textbf{Question 3.} Augmentation de $0{,}5\,\%$ : $CM = 1 + \dfrac{0{,}5}{100} = 1{,}005$.

\textbf{Question 4.} Réduction de $100\,\%$ : $CM = 1 - \dfrac{100}{100} = 0$. La grandeur est annulée.

\end{corrige}

%---------------------------------------------------------------

\begin{exercice}{1SPE-SUITES-FR-R2-EX2}{1}{6}
Un capital de $2\,000$~€ est placé sur un livret qui rapporte $3\,\%$ d'intérêts par an.
\begin{enumerate}
  \item Quel est le coefficient multiplicateur annuel ?
  \item Quel est le capital après $1$~an ? Après $2$~ans ?
  \item Exprimer le capital $C_n$ après $n$ années en fonction de $n$.
  \item Calculer le capital après $10$~ans (arrondir à l'euro).
\end{enumerate}
\end{exercice}

% BEGIN-VERIFY
% # Capital apres n annees
% C0 = 2000
% q = 1.03
% C1 = C0 * q
% assert abs(C1 - 2060) < 1e-6
% C2 = C0 * q**2
% assert abs(C2 - 2121.80) < 0.01
% C10 = C0 * q**10
% assert abs(C10 - 2687.0) < 1.0
% END-VERIFY

\begin{corrige}{1SPE-SUITES-FR-R2-EX2}

\commentaireMarge{$+3\,\%$ par an correspond à multiplier par $1{,}03$ chaque année.}

\textbf{Question 1.} $CM = 1 + \dfrac{3}{100} = 1{,}03$.

\textbf{Question 2.}
\[
C_1 = 2\,000 \times 1{,}03 = 2\,060 \text{ €.}
\]
\[
C_2 = C_1 \times 1{,}03 = 2\,060 \times 1{,}03 = 2\,121{,}80 \text{ €.}
\]

\commentaireMarge{On applique $n$ fois le CM à partir de $C_0$ : formule d'une suite géométrique.}

\textbf{Question 3.} À chaque année, le capital est multiplié par $1{,}03$. Donc :
\[
C_n = 2\,000 \times 1{,}03^n.
\]
La suite $(C_n)$ est une suite géométrique de premier terme $C_0 = 2\,000$ et de raison $q = 1{,}03$.

\textbf{Question 4.}
\[
C_{10} = 2\,000 \times 1{,}03^{10} \approx 2\,000 \times 1{,}3439 \approx 2\,688 \text{ €.}
\]

\end{corrige}

%---------------------------------------------------------------

\begin{exercice}{1SPE-SUITES-FR-R2-EX3}{1}{6}
Un article vaut $500$~€. Son prix augmente de $10\,\%$ la première année, puis diminue de $10\,\%$ la deuxième année.
\begin{enumerate}
  \item Calculer le prix après la première année.
  \item Calculer le prix après la deuxième année.
  \item Le prix est-il revenu à $500$~€ ? Justifier à l'aide des coefficients multiplicateurs.
  \item Quel taux global d'évolution (en \%) cela représente-t-il sur les deux ans ?
\end{enumerate}
\end{exercice}

% BEGIN-VERIFY
% P0 = 500
% P1 = P0 * 1.10
% assert abs(P1 - 550) < 1e-9
% P2 = P1 * 0.90
% assert abs(P2 - 495) < 1e-9
% CM_global = 1.10 * 0.90
% assert abs(CM_global - 0.99) < 1e-9
% taux_global = (CM_global - 1) * 100
% assert abs(taux_global - (-1)) < 1e-9
% END-VERIFY

\begin{corrige}{1SPE-SUITES-FR-R2-EX3}

\textbf{Question 1.}
\[
P_1 = 500 \times 1{,}10 = 550 \text{ €.}
\]

\textbf{Question 2.}
\[
P_2 = 550 \times 0{,}90 = 495 \text{ €.}
\]

\commentaireMarge{Les évolutions successives se composent en multipliant les CM, pas en additionnant les taux.}

\textbf{Question 3.} Non. Le coefficient multiplicateur global est $1{,}10 \times 0{,}90 = 0{,}99 \neq 1$. Une hausse de $10\,\%$ suivie d'une baisse de $10\,\%$ ne se compensent pas exactement.

\textbf{Question 4.} $CM_{global} = 0{,}99$, donc le taux global est :
\[
t = (0{,}99 - 1) \times 100 = -1\,\%.
\]
Sur deux ans, le prix a diminué de $1\,\%$.

\end{corrige}

%---------------------------------------------------------------

\begin{exercice}{1SPE-SUITES-FR-R2-EX4}{1}{7}
Un bactérie se multiplie : la population double toutes les heures. Au départ, on compte $500$ bactéries.
\begin{enumerate}
  \item Quel est le taux d'augmentation horaire ? Le coefficient multiplicateur ?
  \item Donner l'expression de la population $P_n$ après $n$ heures.
  \item Après combien d'heures la population dépasse-t-elle $10\,000$ bactéries ?\\
  (On pourra essayer les valeurs $n = 1, 2, 3, 4, 5$.)
  \item Quel serait le coefficient multiplicateur si la population triplait chaque heure ?
\end{enumerate}
\end{exercice}

% BEGIN-VERIFY
% P0 = 500
% q = 2
% # Population apres n heures
% Pn = lambda n: P0 * q**n
% assert Pn(0) == 500
% assert Pn(1) == 1000
% assert Pn(4) == 8000
% assert Pn(5) == 16000
% # Seuil 10000 : n=5 car P4=8000 < 10000 <= 16000=P5
% assert Pn(4) < 10000 and Pn(5) >= 10000
% # Tripler : CM = 3
% CM_triple = 3
% assert CM_triple == 3
% END-VERIFY

\begin{corrige}{1SPE-SUITES-FR-R2-EX4}

\commentaireMarge{Doubler signifie multiplier par $2$, soit une augmentation de $100\,\%$.}

\textbf{Question 1.} La population double : le taux est $+100\,\%$ et le coefficient multiplicateur est $CM = 2$.

\textbf{Question 2.}
\[
P_n = 500 \times 2^n.
\]

\textbf{Question 3.} On calcule les premières valeurs :

\begin{center}
\begin{tabular}{|c|c|}
\hline
$n$ & $P_n$ \\
\hline
$0$ & $500$ \\
$1$ & $1\,000$ \\
$2$ & $2\,000$ \\
$3$ & $4\,000$ \\
$4$ & $8\,000$ \\
$5$ & $16\,000$ \\
\hline
\end{tabular}
\end{center}

La population dépasse $10\,000$ à partir de $n = 5$, c'est-à-dire après $5$~heures.

\textbf{Question 4.} Si la population triplait, $CM = 3$ (augmentation de $200\,\%$).

\end{corrige}
